%% ============================================================================
%%  applications/template.tex  —  SKELETON RESUME (single source template)
%% ----------------------------------------------------------------------------
%%  This is the ONLY resume template. The orchestrator copies this file into a
%%  position folder, then the writer fills the Experience / Projects / Skills
%%  sections from the context.md bullet pool. It replaces the old per-track
%%  templates (full-stack / ai-ml / dev-ops) — track is now a doctrine-routing
%%  label, not a file choice, so any track (including robotics/autonomy) fills
%%  this same skeleton.
%%
%%  WHAT IS FIXED (do not change): header (\name, \contact) and Education.
%%  WHAT GETS FILLED: Experience entries, Projects entries, Technical Skills.
%%
%%  rSectionEntry signature — 5 positional params, ALWAYS supply all 5 braces:
%%     {1 name}{2 dates}{3 tagline/role/link}{4 technologies}{5 location}
%%     - param 4 (technologies) renders after the name with a " | " separator;
%%       leave it {} for experiences, fill it for projects (the tech-stack line).
%%     - any unused param is an empty {}.
%%  rSet signature — {1 bucket label}{2 comma-separated items}.
%%
%%  PATH NOTE: this template sits at resume/applications/template.tex, beside
%%  template.cls, so the class path is just `template`. On copy into
%%  applications/{year}/{company}/{role}/, the orchestrator rewrites it to
%%  ../../../template (three levels up to applications/, where template.cls lives).
%% ============================================================================

\documentclass[
	11pt
]{../../../template}

\name{Ankur Desai}
\contact{(248) $\hspace{-0.1em}$ 657 $\hspace{-0.4em}$ -- $\hspace{-0.4em}$ 3805 \\ ardusa05@gmail.com \\ \href{https://ardusa.dev/}{www.ardusa.dev} \\ \href{https://linkedin.com/in/ardusa/}{linkedin.com/in/ardusa}}

\begin{document}

	% ---- EDUCATION (fixed; writer does not modify) --------------------------
	\begin{rSection}{E}{ducation}
		\begin{rSectionEntry}
			{Michigan State University}
			{Expected May 2028}
			{B.S. in Computer Science}
			{}
			{East Lansing, MI}

			\item \textbf{GPA:} 3.5 / 4.0
			\item \textbf{Coursework:} Data Structures \& Algorithms, OOP, Computer Architecture, Linear Algebra, Discrete Math
		\end{rSectionEntry}
	\end{rSection}

	\begin{rSection}{E}{xperience}
		\begin{rSectionEntry}
			{Secure and Efficient Autonomous Systems Lab}
			{May 2026 -- Present}
			{Undergraduate Research Assistant}
			{}
			{East Lansing, MI}

			\item Stress-tested the security of an LLM-controlled robot, targeting the cheap per-step gate that screens its motion plans as the exact seam an attacker slips a harmful plan through, one innocent fragment at a time.
			\item Reproduced a published adversarial-robotics attack loop (attacker, target, judge, syntax-checker) as the baseline, then drove the decomposition attack to an 80\% success rate against the per-step triage gate.
		\end{rSectionEntry}

		\begin{rSectionEntry}
			{MindMosaic}
			{Aug 2025 -- May 2026}
			{Software Engineering Intern}
			{}
			{East Lansing, MI}

			\item Built and shipped an interactive knowledge-graph platform that keeps dense, interconnected content navigable as it grows, delivering a production MVP to real users on a polyglot Java Spring Boot, Neo4j, and PostgreSQL stack.
			\item Migrated hot read paths off recursive SQL self-joins onto a Neo4j graph layer so relationships stay first-class, cutting query latency 48\% (850ms to 440ms) and making multi-hop traversals 3x faster than the equivalent SQL.
			\item Designed a dual-database layer that keeps transactional records in PostgreSQL while serving relationship-heavy reads from Neo4j, routing each query class to the store that answers it fastest instead of forcing one engine to do both jobs.
		\end{rSectionEntry}

	\end{rSection}

	\begin{rSection}{P}{rojects}
		\begin{rSectionEntry}
			{Dadei}
			{\href{https://dadei.app}{dadei.app}}
			{Ambient voice assistant that turns overheard conversation into human-approved calendar, email, and task actions}
			{Python, FastAPI, Redis Streams, pgvector, React, Docker}
			{}

			\item Routed all model-proposed side effects through one approval chokepoint with a cancelable countdown, giving every irreversible action one auditable seam between the LLM and the outside world.
			\item Engineered a no-merge-first speaker-ID model using an EMA voice centroid plus a bounded prototype bank with floor-and-margin gating, abstaining into a new identity on ambiguity, attributing 96\% of utterances to the correct speaker.
			\item Replaced single-signal context lookup with hybrid retrieval over pgvector that fuses vector similarity, lexical overlap, recency decay, and participant overlap, grounding each proposal under a fixed token budget.
			\item Decoupled background workers from the live FastAPI service behind a Redis Streams event bus with consumer groups, dead-letter routing after 5 retries, and stale-claim recovery, giving at-least-once delivery that pub/sub dropped.
		\end{rSectionEntry}

		\begin{rSectionEntry}
			{fliks}
			{\href{https://github.com/fliks-gg}{github.com/fliks-gg}}
			{Multi-game platform that crowd-validates gameplay skill and mints shareable certification cards}
			{Go, PostgreSQL, Object Storage, ML Rating Service}
			{}

			\item Built a self-hosted video transcoding pipeline with a concurrent Go worker that pulls upload jobs, transcodes, and writes processed video back to object storage, owning the media path end to end instead of renting a managed service.
			\item Architected an ML rating service that grades each clip independently and shows its verdict beside the crowd's, with those crowdsourced ratings serving as the labeled data the classifier trains on.
		\end{rSectionEntry}

		\begin{rSectionEntry}
			{WizViz}
			{\href{https://devpost.com/software/wizviz}{devpost.com/software/wizviz}}
			{Webcam fighting game driven by real-time body-gesture recognition, winner of the Interactive Media track at SpartaHack X}
			{Python, MediaPipe, OpenCV}
			{}

			\item Classified five distinct combat gestures from body landmarks with a scale-invariant geometric model normalized by torso length, making detection work across players and camera distances without any labeled training data.
		\end{rSectionEntry}
	\end{rSection}

	\begin{rSection}{T}{echnical Skills}
		\begin{rSet}{Languages}{Python, Java, Go, JavaScript/TypeScript, C++, SQL}
		\end{rSet}
		\begin{rSet}{AI / ML}{Large Language Models, Machine Learning, Agentic Workflows, OpenCV}
		\end{rSet}
		\begin{rSet}{Databases}{PostgreSQL, Neo4j, Redis, pgvector}
		\end{rSet}
		\begin{rSet}{Cloud \& Infrastructure}{AWS (EC2, S3, RDS), Docker}
		\end{rSet}
		\begin{rSet}{Tools}{Git}
		\end{rSet}
	\end{rSection}

\end{document}